\section{The exported theorems and the classical reductions}\label{sec:top}

\subsection{What Main.lean states}

\texttt{Main.lean} imports two files: \texttt{Statements.lean}, which imports only Mathlib, and
\nolinkurl{Supplements/Kernels/MainTheoremsAll.lean} \stlx{src/Main.lean}{1}. Its first theorem has
Mathlib's proposition as its type:
\leanquote{riemannHypothesis}
\noindent Mathlib defines \mathlib{RiemannHypothesis} as: every $s\in\mathbb C$ with $\zeta(s)=0$,
$s\ne-2(n+1)$ for all $n\in\mathbb N$, and $s\ne1$ has $\operatorname{Re}s=\tfrac12$
\stc{mathlib2020}. (The definition explicitly excludes the pole $s=1$. Mathlib's
\mathlib{riemannZeta} is a total function and assigns a finite non-zero value there.) The tree restates this proposition in its
own words, and the second exported theorem proves the restatement with the same proof term, so the
two propositions must be definitionally equal for the file to compile:
\leanquote{RiemannZetaZerosOnLine}
\noindent The other three exported theorems are the generalized Riemann hypothesis for Dirichlet
$L$-functions in three forms: non-vanishing on $\operatorname{Re}s>\frac12$ (\lean{LFunction_ne_zero_of_half_lt_re},
\leanloc{src/Main.lean}{8}), zeros in the critical strip lie on the line for every character
(\leanloc{src/Main.lean}{10}) and for primitive characters (\leanloc{src/Main.lean}{12}). The first
of these is the statement everything else reduces to:
\leanquote{DirichletLNonvanishingRightOfHalf}

\subsection{One proposition carries the whole proof}

The four theorems in \texttt{Supplements/Kernels/MainTheorems.lean} that \texttt{Main.lean} re-exports
all take a single input, \lean{allChi}:
\leanquote{allChi}
\leanquote{rhOfAll}
\noindent where \lean{UpperTheory.EveryCharacter} is, word for word, the proposition
\lean{DirichletLNonvanishingRightOfHalf} again:
\leanquote{UpperTheory.EveryCharacter}

\begin{statement}[Every character]\label{st:every}
For every $N\ge1$, every Dirichlet character $\chi$ modulo $N$ and every $s$ with
$\operatorname{Re}s>\frac12$, except $\chi=1$ together with $s=1$, one has $L(s,\chi)\neq0$.
\stl{src/Development/AutomorphicArithmetic/CharacterBoxes/EveryCharacter.lean}{13}
\end{statement}

The rest of the traced source supplies the proof body of Statement~\ref{st:every}, the theorem
\lean{UpperTheory.FinalTheorems.everyCharacter} (\S\ref{sec:architecture}). The reductions from it
to the five exported theorems are short and use only Mathlib.

\subsection{The reductions}

\begin{enumerate}
\item \emph{Riemann hypothesis.} \lean{Formalism.VW.rh_of_every} specializes
  Statement~\ref{st:every} to the trivial character modulo $1$, whose $L$-function is $\zeta$
  (Mathlib's \mathlib{LFunction_modOne_eq}), and passes to \lean{Formalism.VW.rh_of_right}:
  \leanquote{Formalism.VW.rh_of_right}
  For a non-trivial zero $\rho$ with $\operatorname{Re}\rho<\frac12$, the proof first excludes
  $\rho=0$ using $\zeta(0)=-\frac12$. It then evaluates the functional equation
  $\zeta(1-s)=2(2\pi)^{-s}\Gamma(s)\cos(\pi s/2)\zeta(s)$ (Mathlib's \mathlib{riemannZeta_one_sub})
  at $s=1-\rho$, where the factors $2(2\pi)^{-s}$, $\Gamma(s)$ and $\cos(\pi s/2)$ are non-zero: a
  vanishing cosine in this range would force $\rho=0$ or a negative even integer
  \stl{src/Formalism/GammaFactorNonvanishing/RhOfRight.lean}{5} \stcs~\cite{davenport}.
\item \emph{Primitive characters, critical strip.} \lean{Formalism.VW.critLine}
  \stl{src/Formalism/GammaFactorNonvanishing/CritLine.lean}{11} does the same with the functional
  equation of the completed $L$-function of a primitive character, which relates $\chi$ at $s$ to
  $\chi^{-1}$ at $1-s$ (Mathlib's \mathlib{IsPrimitive.completedLFunction_one_sub}); the
  archimedean factor $\Gamma_{\mathbb R}(s+\kappa)$, $\kappa\in\{0,1\}$, is non-zero for
  $\operatorname{Re}s>0$ (\lean{Formalism.VW.gammaFactor_ne_zero},
  \leanloc{src/Formalism/GammaFactorNonvanishing/CritLine.lean}{6}).
\item \emph{All characters, critical strip.} \lean{Claim.strip}
  \stl{src/Development/AutomorphicArithmetic/CharacterBoxes/ClaimStrip.lean}{59} first passes from
  $\chi$ to the primitive character $\chi^\ast$ inducing it: in $0<\operatorname{Re}s<1$,
  $L(s,\chi)=L(s,\chi^\ast)\prod_{p\mid N}(1-\chi^\ast(p)p^{-s})$ and the Euler factors do not vanish
  because $|\chi^\ast(p)p^{-s}|<1$ (\lean{Claim.prim_zero}, \lean{Claim.factor_ne_zero}).
\item \emph{The statement for $\zeta$ in Mathlib's form} also follows directly (\lean{Claim.zeta},
  \leanloc{src/Development/AutomorphicArithmetic/CharacterBoxes/ClaimStrip.lean}{75}), using
  $\zeta(s)\ne0$ for $\operatorname{Re}s\ge1$ (Mathlib's \mathlib{riemannZeta_ne_zero_of_one_le_re}).
\end{enumerate}

These steps are the textbook passage from ``no zeros in $\operatorname{Re}s>\frac12$'' to ``all
non-trivial zeros on the line'' \stcs~\cite[Ch.~8--9]{davenport}. A second, parallel reduction,
\lean{UpperTheory.riemannHypothesis_of_everyCharacter}
\stl{src/Development/AutomorphicArithmetic/CharacterBoxes/RiemannHypothesisOfEveryCharacter.lean}{215},
is used by \lean{UpperTheory.FinalTheorems.riemannHypothesis} but not by \texttt{Main.lean}.
The kernel files \texttt{MainTheoremsAll.lean} and \texttt{MainTheoremsCritical.lean} (two and four
lines) state no theorem; they only import four further files on small class numbers.
